\subsection{معرفی و انگیزه}
این مکانیزم بر مبنای مبانی نظری پیشرفته در ادبیات توجه و فشرده‌سازی \cite{ref_jamba_2024} توسعه یافته است.

معماری \lr{Jamba} یک الگوی ساختاری پیشگام و ترکیبی (\lr{Hybrid Transformer-SSM}) است که با هدف ادغام مزایای بنیادین مدل‌های فضای حالت گزینشی (\lr{Selective State Space Models}) نظیر \lr{Mamba} با توانایی بازیابی دقیق و همبستگی دوربرد مکانیزم توجه چندسر (\lr{Multi-Head Attention - MHA}) به همراه افزایش ظرفیت پارامتری از طریق لایه‌های خبره ناهمگن (\lr{Expert Choice Mixture of Experts - MoE}) معرفی گردیده است. 

در مدل‌های استاندارد توجه چندسر، هزینه محاسباتی و اندازه کش حافظه برای ذخیره تانسورهای کلید و مقدار (\lr{KV-Cache}) به صورت خطی با طول زمینه رشد می‌کند که عملکرد مدل در زمینه‌های طویل را به شدت محدود می‌نماید. در نقطه مقابل، اگرچه مدل‌های فضای حالت با وضعیت پنهان متراکم خود دارای هزینه زمانی $\mathcal{O}(1)$ در مرحله تولید توکن و پیچیدگی خطی در طول توالی هستند، اما در وظایفی که نیازمند توجه دقیق به تمام نشانه‌های قبلی در پنجره‌های متن حجیم است (همچون استدلال و پاسخگویی به سوالات طولانی)، با افت ظرفیت مواجه می‌گردند. ساختار \lr{Jamba} با درهم‌آمیزی متناوب شاخه‌های \lr{SVD-Mamba}، لایه سنتی \lr{MHA} با سازوکارهای مدرن \lr{GQA}، و سامانه گزینش هوشمند خبرگان (\lr{Expert Choice MoE})، پاسخی یکپارچه به چالش مقیاس‌پذیری طول دنباله بدون فدا کردن دقت ارائه می‌دهد.
